\documentclass[10pt,a4paper]{amsart}
\usepackage[margin=19mm]{geometry}
\usepackage{amsmath,amssymb,mathtools}
\usepackage[scr=rsfs,cal=euler]{mathalfa}
\usepackage{xfrac}
\usepackage[unicode,hidelinks,bookmarks=false]{hyperref}
\newcommand{\ie}{i.e.\ }
\newcommand{\cf}{cf.\ }
\newcommand{\resp}{respectively\ }
\newcommand{\Subsection}{\subsection}
\providecommand{\nobreakdash}{\nobreak\hbox{-}}
\setlength{\emergencystretch}{2em}
\multlinegap0pt
\usepackage{bm}
\usepackage{bbm}
\usepackage{dsfont}
\usepackage{xspace}
\usepackage{cancel}
\usepackage{mathtools}
\usepackage{amsthm}
\makeatletter
\@ifundefined{prop}{\newtheorem{prop}{Proposition}}{}
\@ifundefined{thm}{\newtheorem{thm}{Theorem}}{}
\makeatother
\DeclareMathOperator*{\argmin}{arg\,min}
\def\Atop {\matr{A}^{\top}}
\def\R {\mathbb{R}}
\def\Rm {\R^m}
\def\bb {\boldsymbol{b}}
\def\bd {\boldsymbol{d}}
\def\be {\boldsymbol{e}}
\def\bp {\boldsymbol{p}}
\def\bpa {\bp_{a}}
\def\bpo {\bp_{0}}
\def\dom {\mathrm{dom~}}    % Domain
\def\epo {\mathcal{E}(\bp_{0})}
\def\proj {\mathrm{proj}}
\def\sign {\mathrm{sgn}}
\newcommand{\matr}[1]{\bm{#1}}
\newcommand{\normsq}[1]{\left\Vert{#1} \right\Vert_{2}^{2}}
\newcommand{\normtwo}[1]{\left\Vert{#1} \right\Vert_{2}}
\title[A fast algorithm for solving the lasso problem exactly witho]{A fast algorithm for solving the lasso problem exactly without homotopy using differential inclusions (Theorem 4.1)}
\author{Gabriel P. Langlois, J\'er\^ome Darbon}
\date{Source: arXiv:2507.05562v2 (CC BY 4.0)}
\begin{document}
\maketitle

\noindent\emph{Source and licence.} A fast algorithm for solving the lasso problem exactly without homotopy using differential inclusions. Version arXiv:2507.05562v2 (CC BY 4.0), distributed under the Creative Commons Attribution 4.0 International licence, as recorded in the arXiv licence metadata for this exact version and shown on the abstract page https://arxiv.org/abs/2507.05562v2. Full rights notice: licenses/arxiv-2507.05562-CC-BY-4.0.txt.\par\smallskip

\noindent\textit{Editorial scope.} Counted unit: the Thm printed as Theorem 4.1 of the article (source0.tex lines 578--603, source label \texttt{thm:mr\_local\_sol}), delivered together with the article's own complete original proof of it (source0.tex lines 605--630, one \texttt{proof} environment of 26 lines). No mathematics is written, repaired or re-derived: statement and proof are the article's own text line for line. Required context retained uncounted: thm:eu\_diff\_inclusions (lines 424--442); eq:general\_diff\_inclusion (lines 425--427); eq:bp\_slow\_ode (lines 458--460); prop:descent\_direction (lines 526--564); eq:identity\_dd (lines 560--562). Every definition, display and label of those lines is retained unchanged. Omitted from this package and not redistributed: 1--423, 428--457, 461--525, 563--577, 604--604, 631--1663. Nothing inside a retained excerpt is omitted or altered: every retained body is delivered exactly as the article's lines are written, so no omission receipt of any kind accompanies this package. Citations: the retained excerpts carry no citation command at all, so no bibliography entry of the article is transcribed and no reference is redistributed. Reference adaptation: none was needed and none was made; every reference inside the delivered unit resolves inside it, so no unresolved reference or citation is produced. Printed slips kept literal: no printed word, symbol, hypothesis or number of the counted statement or of its proof was altered. Layout and class: the article's own class is \texttt{article} with packages including geometry, bm, xcolor, setspace, natbib, appendix, enumerate, algorithm2e, multirow, graphicx, subcaption, tabularx, makecell, booktabs; the wrapper is the lane's pinned \texttt{amsart} editorial wrapper at 10pt on A4, which restates the theorem-like environments the retained text needs and adds the article's own macro definitions that the retained text needs (\texttt{\textbackslash Atop}, \texttt{\textbackslash R}, \texttt{\textbackslash Rm}, \texttt{\textbackslash argmin}, \texttt{\textbackslash bb}, \texttt{\textbackslash bd}, \texttt{\textbackslash be}, \texttt{\textbackslash bp}, \texttt{\textbackslash bpa}, \texttt{\textbackslash bpo}, \texttt{\textbackslash dom}, \texttt{\textbackslash epo}, \texttt{\textbackslash matr}, \texttt{\textbackslash normsq}, \texttt{\textbackslash normtwo}, \texttt{\textbackslash proj}, \texttt{\textbackslash sign}), with extra wrapper packages (bm, bbm, dsfont, xspace, cancel, mathtools). The delivered numbering is the unit's own. Assets: no retained span contains a figure environment, a graphics-inclusion command, a PDF-page inclusion, a table or a TikZ picture; no image file is copied into this package, the package declares no asset dependency and no image file is redistributed with it. Licence: version arXiv:2507.05562v2 of this article is distributed under the Creative Commons Attribution 4.0 International licence, as recorded in the arXiv licence metadata for this exact version and shown on the abstract page https://arxiv.org/abs/2507.05562v2. A full rights notice is delivered at licenses/arxiv-2507.05562-CC-BY-4.0.txt. Characters: every retained excerpt is pure ASCII, so the delivered engine, XeLaTeX with its default Latin Modern text font, reports no missing character.
\begin{thm}[Existence and uniqueness of solutions to gradient inclusions]\label{thm:eu_diff_inclusions} Let $g\colon \Rm \to \R\cup\{+\infty\}$ be a proper, lower semicontinuous, and convex function and let $\bp_{0} \in \dom \partial g$. Consider the gradient inclusions
\begin{equation}\label{eq:general_diff_inclusion}
    \dot{\bp}(\tau) \in - \partial g(\bp(\tau)), \qquad \bp(0) = \bp_{0}.
\end{equation}
Then, there exists a unique solution $\bp \colon [0,+\infty) \mapsto \dom \partial g$ satisfying~\eqref{eq:general_diff_inclusion}. Moreover:
\begin{itemize}
    \item[(i)] The function $\tau \mapsto \dot{\bp}(\tau)$ is right-continuous and the function $\tau \mapsto \normtwo{\dot{\bp}(\tau)}$ is nonincreasing.
    
    \item[(ii)] If $g$ achieves its minimum at some point, then $\bp(\cdot)$ converges to such a point:
    \[
    \lim_{\tau \to +\infty} g(\bp(\tau)) = \min_{\bp \in \Rm} g(\bp) \quad \text{and} \quad \lim_{\tau \to +\infty} \bp(\tau) \in \argmin_{\bp \in \Rm} g(\bp).
    \]
    \item[(iii)] (The minimal selection principle.) The function $\bp(\cdot)$ satisfies the initial value problem
    \begin{equation}\label{eq:general_slow_ode}
            \dot{\bp}(\tau) =-\proj_{\partial g(\bp(\tau))}(\boldsymbol{0}), \qquad \bp(0) = \bp_{0}
    \end{equation}
    at $\tau = 0$ and almost everywhere on $(0,+\infty)$. 
\end{itemize}
\end{thm}

\begin{equation}\label{eq:bp_slow_ode}
    \dot{\bp}(\tau) =-\proj_{\partial_{\bp} V(\bp(\tau);t,\bb)}(\boldsymbol{0}), \qquad \bp(0) \in \dom \partial_{\bp} V(\cdot;t,\bb)
\end{equation}

\begin{prop}\label{prop:descent_direction}
Let $t \geqslant 0$ and $\bpo \in \dom \partial_{\bp} V(\cdot;t,\bb)$. Then:
    \begin{itemize}
    \item[(i)] There exists $\Delta_{*}(\bpo;t,\bb) > 0$ such that $\bpo + \Delta\bd(\bpo;t,\bb) \in \dom V(\cdot;t,\bb)$ for every $\Delta \in [0,\Delta_{*}(\bpo;t,\bb)]$, where
    \begin{equation}\label{eq:descent_max_time}
    \Delta_{*}(\bpo;t,\bb) = \min_{j \in \{1,\dots,n\}} \left\{\frac{\sign{\left\langle \matr{D}(\bpo)\Atop\bd(\bpo;t,\bb),\be_j\right\rangle} - \left\langle \matr{D}(\bpo)\Atop\bpo,\be_j\right\rangle)}{\left\langle\matr{D}(\bpo) \Atop\bd(\bpo;t,\bb),\be_j\right\rangle} \right\}.
    \end{equation}
    Moreover, $\Delta_{*}(\bpo;t,\bb) = +\infty \iff \Atop\bd(\bpo;t,\bb) = \boldsymbol{0} \iff \bd(\bpo;t,\bb) = \boldsymbol{0}$.


    \item[(ii)](Descent direction) For every $\Delta \in [0, \Delta_{*}(\bpo;t,\bb)]$, we have
    \begin{equation}\label{eq:V_formula_descent}
    V(\bpo + \Delta\bd(\bpo;t,\bb);t,\bb) - V(\bpo;t,\bb) = -\Delta\left(1-t\Delta/2\right)\normsq{\bd(\bpo;t,\bb)}.
    \end{equation}
    In particular, if $\bd(\bpo;t,\bb) \neq \boldsymbol{0}$, then it is a descent direction of $V(\cdot;t,\bb)$. 


    \item[(iii)] For every $\Delta \in [0, \Delta_{*}(\bpo;t,\bb)]$, we have the inclusion
    \[
    -(1-t \Delta)\bd(\bpo;t,\bb) \in \partial_{\bp} V(\bpo + \Delta\bd(\bpo;t,\bb);t,\bb).
    \]


    \item[(iv)] For every $\Delta \in [0,\Delta_{*}(\bpo;t,\bb))$, we have the inclusions
    \[
    \mathcal{E}(\bpo + \Delta\bd(\bpo;t,\bb)) \subset \epo
    \]
    and
    \[
    \partial_{\bp}V(\bpo + \Delta\bd(\bpo;t,\bb);t,\bb) \subset \{t\Delta\bd(\bpo;t,\bb)\} + \partial_{\bp}V(\bpo;t,\bb).
    \]


    \item[(v)] (Evolution rule) For every $\Delta \in [0,\Delta_{*}(\bpo;t,\bb))$, the following evolution rule holds:
    \begin{equation}\label{eq:identity_dd}
    \bd(\bpo + \Delta\bd(\bpo;t,\bb);t,\bb) = (1-t \Delta)\bd(\bpo;t,\bb).
    \end{equation}
\end{itemize}
\end{prop}

\begin{thm}\label{thm:mr_local_sol}
Let $t \geqslant 0$ and $\bpo \in \dom \partial_{\bp} V(\cdot;t,\bb)$. The slow system~\eqref{eq:bp_slow_ode} with initial value $\bp(0) = \bpo$ coincides with the initial value problem
\begin{equation}\label{eq:dd_decay}
\dot{\bp}(\tau) = e^{-t\tau}\bd(\bpo;t,\bb), \quad \bp(0) = \bpo
\end{equation}
on $\tau \in [0,\tau_{*}(\bpo;t,\bb))$, where 
\begin{equation}\label{eq:thm_tau0}
\tau_{*}(\bpo;t,\bb) = 
\begin{cases}
    \Delta_{*}(\bpo;0,\bb) & \text{if} \ t = 0, \\
    -\ln\left(1-t\Delta_{*}(\bpo;t,\bb)\right)/t & \text{if} \ t > 0 \ \text{and} \ 1-t\Delta_{*}(\bpo;t,\bb) > 0, \\
    +\infty & \text{otherwise}.
\end{cases}
\end{equation}
In particular, the slow solution is given explicitly on $[0,\tau_{*}(\bpo;t,\bb))$ by
\begin{equation}\label{eq:slow_local_solution}
    \bp(\tau) = \bpo + f(\tau,t)\bd(\bpo;t,\bb),
\end{equation}
where 
\begin{equation}\label{eq:f_fun}
f(\tau,t) = \begin{cases}
    \tau & \text{if} \ t = 0,\\
    \left(1 - e^{-t\tau}\right)/t & \text{if} \ t > 0.
\end{cases}
\end{equation}
\end{thm}

\begin{proof}
We will use the evolution rule~\eqref{eq:identity_dd} to show that the affine system
\begin{equation}\label{eq:affine_system}
    \dot{\bp}_{a}(\tau) =-t(\bpa(\tau) - \bpo) + \bd(\bpo;t,\bb), \qquad \bpa(0) = \bpo.
\end{equation}
satisfies the slow system~\eqref{eq:bp_slow_ode} on $[0,\tau_{*}(\bpo;t,\bb))$ and conclude using uniqueness. First, consider the affine system~\eqref{eq:affine_system}. A short calculation shows that its unique, global solution is given by 
\[
\bpa(\tau) = \bpo + f(\tau,t)\bd(\bpo;t,\bb).
\]
Substitute the solution in~\eqref{eq:affine_system} above to find
\begin{equation}\label{eq:explicit_slow_evolution}
\begin{alignedat}{1}
\dot{\bp}_{a}(\tau) &= (1-tf(\tau,t))\bd(\bpo;t,\bb) \\
&= e^{-t\tau}\bd(\bpo;t,\bb).
\end{alignedat}
\end{equation}
Next, we invoke the evolution rule~\eqref{eq:identity_dd} in Proposition~\ref{prop:descent_direction}(v) with $\Delta = f(\tau,t)$ to find
\begin{equation}\label{eq:thm_local_id_holds}
-\proj_{\partial_{\bp} V(\bpo + f(\tau,t)\bd(\bpo;t,\bb),t)}(\boldsymbol{0}) = (1-tf(\tau,t))\bd(\bpo;t,\bb)
\end{equation}
whenever $0 \leqslant f(\tau,t) < \Delta_{*}(\bpo;t,\bb)$. Notice $\tau \mapsto f(\tau,t)$ increases monotonically from $0$ to $1/t$ (0 to $+\infty$) when $t > 0$ ($t = 0$). Hence if $t = 0$, the largest value of $\tau$ for which~\eqref{eq:thm_local_id_holds} holds is $\Delta_{*}(\bpo;0,\bb)$. If $t > 0$ and $1-t\Delta_{*}(\bpo;t,\bb) > 0$, then $(1-e^{-t\tau})/t$ is equal to $\Delta_{*}(\bpo;t,\bb)$ at $\tau = -\ln\left(1-t\Delta_{*}(\bpo;t,\bb)\right)/t$. If $t > 0$ and $1 - t\Delta_{*}(\bpo;t,\bb)) \leqslant 0$, then~\eqref{eq:thm_local_id_holds} holds for every $\tau \geqslant 0$. Taken together, we find 
\[
\dot{\bp}_{a}(\tau) = -\proj_{\partial_{\bp} V(\bpo + f(\tau,t)\bd(\bpo;t,\bb),t)}(\boldsymbol{0}) \ \text{for every $\tau \in [0,\tau_{*}(\bpo;t,\bb))$}.
\]
Uniqueness follows from Theorem~\ref{thm:eu_diff_inclusions}, and hence~\eqref{eq:dd_decay} follows from~\eqref{eq:explicit_slow_evolution}.
\end{proof}

\end{document}
